% Legacy segmented-LT global-order proof (not rendered in layout preview).
\section{To be revised:Global Orderedness of Segmented Loser Tree}
\label{sec:appendix-segment-correctness}

For an ordered sequence $S$, let $\operatorname{Pref}_q(S)$ denote its
first $\min\{q,|S|\}$ items. We distinguish two properties: every emitted
segment must be the next prefix of the global order, and every input item
must be emitted exactly once.

\paragraph{Why path selection matters.}
A locally ordered replay does not by itself guarantee global order.
Consider $h=2$, $L=[1,4]$, and $R=[2,6]$. Their replay produces
\[
tmp=[1,2],\qquad \mathit{Loser}(v)=[4,6],
\]
and the left slot is exhausted first. Suppose its next segment begins
with $5$, whereas the next right segment is $[7,8]$. If the algorithm
incorrectly refreshes the right path, replaying $[4,6]$ with $[7,8]$
would emit $[4,6]$, placing $6$ before the unseen item $5$. The replay
batch is locally sorted, but it omits a candidate required for global
order. Hence correctness depends on proving that
$\mathit{Pointer}(v)$ selects exactly the child whose next segment may
precede part of $\mathit{Loser}(v)$.

A replay segment is associated with a \emph{logical child slot}.
This association identifies the child whose continuation is required
when the segment is exhausted; it does not require every item in the
segment to originate from that child. After a replay, each internal node
maintains the following invariant:

\begin{itemize}
  \item $\mathit{Loser}(v)$ is the ordered segment of exposed items not propagated to the parent;
  \item if $\mathit{Pointer}(v)=p$, the next segment from $p$ is the only unexposed segment that may precede an item in $\mathit{Loser}(v)$; and
  \item if $\mathit{Pointer}(v)=\textsc{null}$, the entire loser segment is certified and may be propagated without another ranking call.
\end{itemize}

\begin{lemma}[First-exhausted replay dependency]
\label{lem:segment-frontier-barrier}
Let $L$ and $R$ be valid $h$-item replay segments associated with the two
logical child slots of $v$, and let
\[
  Z=\operatorname{Sort}_{\rho}(L\uplus R)
\]
under a rank map consistent with $\prec_\phi$. Let $p$ be the slot whose
$h$-th item is encountered first when scanning $Z$, and let $o$ be the
other slot. Then
\[
  tmp=Z[1..h],
  \qquad
  \mathit{Loser}(v)=Z[h+1..2h],
  \qquad
  \mathit{Pointer}(v)=p
\]
satisfy the node invariant. In particular, $tmp$ is the next certified
$h$-item prefix, and refreshing any slot other than $p$ is insufficient
to certify the next winner in general.
\end{lemma}

\begin{proof}
Let $b$ be the position at which $p$ is first exhausted. Since both
inputs contain $h$ items,
\[
  h\le b\le 2h-1.
\]
Consequently, every item of $Z[1..h]$ occurs no later than the last
exposed item of $p$. Any subsequent item from $p$ must follow that last
item, and therefore cannot precede $Z[1..h]$. The other slot $o$ is not
exhausted at position $b$; after $b$, all remaining items of $Z$ come
from $o$. Its unseen continuation therefore follows the final item of
$Z$. Thus $Z[1..h]$ is certified against both continuations.

In contrast, the next item from $p$ may lie between $Z[b]$ and
$Z[b+1]$. It may therefore precede part of
$Z[h+1..2h]$, whereas an unseen item from $o$ cannot. Hence $p$ is the
unique unresolved dependency and must become $\mathit{Pointer}(v)$.

The loser segment may contain items originating from both subtrees, but
its suffix following position $b$ consists only of items from $o$ and
ends with the exposed frontier of $o$. It can therefore occupy the
logical slot associated with $o$ during the next replay. If the new
segment from $p$ is exhausted first, $p$ remains the unresolved
dependency. If the stored loser is exhausted first, the exposed frontier
of $o$ has been consumed and the dependency switches to $o$. This is
exactly the pointer update performed by
Algorithm~\ref{alg:replay_segment_node}.
\end{proof}

Algorithm~\ref{alg:replay_segment_node} presents the steady-state case
of two full segments. For a final short segment, its actual length is
used as the exhaustion threshold. A short segment has no continuation;
therefore, if it is exhausted first, the pointer moves to the other slot
when that slot has more output, and becomes \textsc{null} otherwise.

\begin{theorem}[Conditional global orderedness]
\label{thm:segment-global-order}
Fix a strict total order $\prec_\phi$ with deterministic tie-breaking.
Suppose that:
\begin{itemize}
  \item[(A1)] every input run is finite and sorted by $\prec_\phi$;
  \item[(A2)] every rank map produced during initialization or path
        refresh is the restriction of the same order $\prec_\phi$;
  \item[(A3)] every unexhausted producer eventually returns its next
        segment, with only its final segment allowed to contain fewer
        than $h$ items; and
  \item[(A4)] the tree depth is at most $d$ and $(d+1)h\le m$.
\end{itemize}
Then Algorithm~\ref{alg:segmented_loser_tree_merge} emits exactly the
sorted merge of its input runs. Any consumer that stops after $K$ items
receives the first $K$ items of the same global order.
\end{theorem}

\begin{proof}
Let
\[
  \mathcal S=
  \operatorname{Sort}_{\prec_\phi}
  (R_1\uplus\cdots\uplus R_k)
\]
be the desired global merge.

\paragraph{Initialization.}
At line~\ref{line:slt-init}, each leaf returns the next prefix of its
sorted run. By (A1), this segment is a certified prefix of that leaf.
The shared rank map $\rho_0$ orders all initial leaf segments according
to $\prec_\phi$ by (A2).

We apply structural induction to \textsc{BuildSubtree}. A leaf returns
its certified segment directly. At an internal node, the recursive calls
return the next winner segments of its left and right subtrees. By
Lemma~\ref{lem:segment-frontier-barrier},
\textsc{ReplaySegmentNode} returns their next merged prefix and stores a
loser/pointer pair satisfying the node invariant. Hence the value
returned by every recursive call is the next certified segment of its
subtree. In particular, the initial root $tmp$ is the first segment of
$\mathcal S$. The returned values are transient: only
$\mathit{Loser}(v)$ and $\mathit{Pointer}(v)$ remain at an internal node.

\paragraph{Pointer-guided refresh.}
After a root winner is emitted, let
\[
  \mathcal P=(u_0,u_1,\ldots,u_t),
  \qquad u_0=\textsc{Root}(T),
\]
be the path followed by Algorithm~\ref{alg:repair_segment_path}.
At every internal node $u_i$, the node invariant states that the child
selected by $\mathit{Pointer}(u_i)$ is the only child whose continuation
may precede part of $\mathit{Loser}(u_i)$. Following any other edge can
therefore omit a necessary candidate, as illustrated by the preceding
counterexample. Following the stored pointers reaches exactly the
unresolved frontier.

If $u_t$ is a leaf, \textsc{FillLeaf} supplies its next certified
segment. If it is an internal node with a null pointer, its loser is
already fully certified and can be promoted directly. In either case,
the terminal node supplies a valid initial $tmp$.

The refresh batch is
\[
  B=\{tmp\}\uplus
    \biguplus_{v\in\mathcal P}\{\mathit{Loser}(v)\}.
\]
By (A2), $\rho_B$ gives the order induced by $\prec_\phi$ on every item
participating in the refresh. Replay then proceeds bottom-up. Assume the
current $tmp$ is the next certified segment returned by the pointer
child of a path node $v$. The stored loser occupies the other logical
slot. Lemma~\ref{lem:segment-frontier-barrier} shows that replay returns
the next certified segment of $v$ and establishes the new loser/pointer
invariant at $v$. Induction from $u_t$ to $u_0$ therefore proves that the
returned root $tmp$ is the next globally certified segment. Nodes
outside $\mathcal P$ require no update because their continuations cannot
precede the stored loser at any path node.

\paragraph{Root output.}
Let $\Pi_s$ be the output after $s$ root segments and let $U_s$ be the
items not yet emitted. We maintain
\[
  \Pi_s=\operatorname{Pref}_{|\Pi_s|}(\mathcal S),
  \qquad
  tmp_s=
  \operatorname{Pref}_{|tmp_s|}
  \left(\operatorname{Sort}_{\prec_\phi}(U_s)\right).
  \tag{I}
\]
Initialization establishes (I) for $s=0$. If (I) holds, appending
$tmp_s$ gives
\[
  \Pi_{s+1}
  =\Pi_s\mathbin{\|}tmp_s
  =\operatorname{Pref}_{|\Pi_s|+|tmp_s|}(\mathcal S).
\]
The pointer-guided refresh proved above returns the next prefix of the
new remainder $U_{s+1}$, reestablishing (I). Thus every emitted segment
extends the same global sorted prefix.

\paragraph{Conservation, capacity, and termination.}
A child transfers ownership of every item in its returned $tmp$ to its
parent. Ranking changes only their order. At a replay node,
\[
  tmp\mathbin{\|}\mathit{Loser}(v)
  =\operatorname{Sort}_{\rho}(L\uplus R),
\]
so every replayed occurrence is assigned to exactly one transient winner
or one persistent loser. Promotion and output transfer ownership without
copying. Therefore no item is lost or duplicated.

A refresh batch contains at most one $h$-item terminal segment and one
loser of at most $h$ items from each of at most $d$ path levels. Hence
\begin{equation}
  |B|\le (d+1)h\le m.
  \label{eq:appendix-repair-capacity}
\end{equation}
Every nonempty root $tmp$ emits at least one item. By (A3), an
unexhausted terminal eventually supplies its next segment; when a final
short segment is consumed, the boundary rule removes its exhausted
dependency. Because all runs are finite, the process terminates only
after every source, transient winner, and stored loser is empty.
Conservation then gives $|\Pi|=|\mathcal S|$, while invariant (I) gives
$\Pi=\mathcal S$.

A pipelined consumer that stops after $K$ items merely truncates this
same sequence and therefore receives
$\operatorname{Pref}_K(\mathcal S)$.
\end{proof}

The theorem depends on Assumption~(A2): enforcing a valid permutation
inside each individual ranking call does not by itself guarantee
consistency between different calls.
